\storyheader{Độc giả duy nhất của tôi
}
{Ngày 30}
{
Việc làm tiểu thuyết gia, cứ như là một thử thách lớn nhất của sự kiên
nhẫn và sáng tạo vậy. Đôi khi tôi sẽ không thể nảy ra một ý tưởng mới,
đôi khi thì chỉ có ý tượng chứ không thể biết cách triển khai như thế
nào. Hoặc nếu như ta có cả trí tưởng tượng, sự sáng tạo, ngôn từ và cả
sự kiên trì thì cũng chưa chắc thứ được viết ra có hợp ý với độc giả hay
không.

Những câu chuyện được sinh ra để cho mọi người có thể đọc.

Nếu không ai đọc, thì chúng được sinh ra có nghĩa lý gì?

Tôi nhả lưng ra ghế, trước mắt là cái máy tính đang hiện lên ứng dụng
đánh chữ à tôi hay sử dụng để viết tiểu thuyết và hiện tại nó vẫn đang
trống không. Đúng vậy, tôi chả viết được một chữ nào cả.

Nhưng mà có thể việc này cũng không phải là cái gì đấy quá nghiêm trọng,
vì suy cho cùng thì có lẽ tôi cũng sẽ thất bại thôi.

Tôi là Akemi, một người có thể gọi là tiểu thuyết gia. Cơ mà nếu viết
truyện rồi đăng lên thôi mà chả có tí bản quyền cũng như chả nổi tiếng
gì thì có tính là tiểu thuyết gia không nhỉ? Mà thôi kệ vậy, nói chung
là thế đấy. Tôi là tiểu thuyết gia hay cái gì đấy đại loại vậy, và
truyện do tôi viết cũng chả có mấy người đọc.

Vì để nâng cao danh tiếng thì tôi đã đăng ký để dự một cuộc thi khá là
nổi tiếng cho các tác giả trẻ. Ai muốn tham gia thì cứ gửi mail cùng vời
một truyện ngắn cho ban tổ chức để họ duyệt, nếu được duyệt qua thành
công thì bạn sẽ phải tiếp tục viết thêm một truyện nữa và gửi cho ban tổ
chức. Deađline thì là trước ngày trao giải khoảng một ngày.

Vậy đấy, và tôi thì đã may mắn được duyệt vì được đánh giá là có vốn từ
phong phú. Những tưởng bản thân đã có thể chứng minh thực lực thì tôi đã
phải đối mặt với thực tế và vấn đề của bản thân. Khi mà tôi thì vẫn phải
học tập và khả năng sáng tạo lẫn lên ý tưởng cốt truyện của tôi không
được tốt cho lắm.

``Kiểu này chắc không kịp deadline rồi.'' Tôi nằm dài lên bàn và than
thở. ``Hay thôi bỏ cuộc luôn cho rồi nhỉ?''

Dù sao thì chắc có lẽ cũng sẽ không thể dành chiến thắng đâu. Trong cuộc
thi thì người ta chỉ quan tâm đến những người đứng đầu, nếu không thể
được xếp vào trong số đó thì chả có tác dụng gì cả.

Đột nhiên, đầu tôi nhói lên một cơn đau như thể bị ai đó đánh vào vậy.

Mà thật ra đúng là có người đánh tôi thật, thậm chí tôi còn biết đó là
ai nữa cơ.

``Đau lắm đấy...'' Tôi ôm lấy đầu mìn và nhìn về phía sau. ``Kouhai mà
đánh senpai như thế là không tốt đâu.''

``Ai là kouhai cơ?'' cô bạn trước mặt ngồi xuống cái ghế bên cạnh tôi.
``Đừng có nghĩ cậu vào câu lạc bộ trước thì là senpai của tớ, chúng ta
bằng tuổi cả đấy.''

``Thì vẫn là senpai còn gì?''

Cậu ấy là Yuka, bạn thân kiêm bạn thuở nhỏ kiêm fan của tôi. Tôi và cậu
ấy chơi chung với nhau từ tiểu học tới giờ và ở độ tuổi đó thì khả năng
viết truyện của tôi vẫn còn là một thứ gì đấy oách xà lách nên cậu ấy mê
như điếu đổ, thành ra đến tận bây giờ thì vẫn thích truyện của tôi dù
cho chúng cũng chả hay ho gì cho lắm.

Yuka cũng là người đã khuyên tôi đăng ký tham dự cuộc thi kia. Ban đầu
là định từ chối rồi đấy nhưng thấy mặt cậu ấy trong tội quá với lại có
tham gia thì cũng chả mất gì nên là tôi cũng xuôi theo.

Mà không mất gì chỉ là tôi nghĩ vậy thôi, chứ tham gia rồi thì tôi đã
mất nhiều thứ lắm luôn ấy chứ. Vừa mất thời gian và vừa mất sức, thời
gian là tiền bạc nên tín ra thì tôi mất toàn bộ tài sản rồi ấy chứ.

``Nhưng sao cậu lại cốc đầu tớ?'' Tôi quay lại với việc viết truyện và
tiện thể hỏi Yuka về lý do cho vết thương trên đầu mình. ``Tớ làm cậu
giận à?''

``Tất nhiên là giận rồi, tại cậu muốn bỏ cuộc chứ gì nữa. Chẳng thèm cố
gắng gì cả.''

``Tớ có mà?'' Tôi quay sang nói với cậu ấy bằng vẻ mặt buồn bã giả tạo
của mình.

``Không, không hề.''

Lần này thì tôi cũng chẳng buồn cãi lại nữa bởi vì cậu ấy nói không sai.
Đúng là tôi có bỏ ra rất nhiều thời gian cho việc ngồi viết mấy cái
truyện vô nghĩa và sẽ xóa ngay sau đó vì không thấy hài lòng, nhưng thực
chất thì việc đó cũng chả có tác dụng lắm. Nói thẳng là thì nó là nỗ lực
ảo ấy.

Tôi chả bao giờ nghiêm túc ngồi nghĩ một cái cốt chuyện lẫn tình tiết
sao cho hoàn toàn nghiêm túc cả. Tất cả những gì tôi làm chỉ là ngồi đốt
thời gian một cách vô ích và xóa mấy nghìn từ mình vừa viết như là tôi
đang làm hiện tại đây.

Tôi có thể làm gì để cải thiện không nhỉ?

``Nhưng dù sao thì cũng đâu có được giải đâu chứ, cuộc thi đó biết bao
nhiêu tác giả trẻ nổi tiếng tham gia mà?'' Tôi xóa xong cái truyện vừa
viết của mình rồi quay sang cậu ấy. ``Tớ thì làm sao mà đủ sức cạnh
tranh với họ?''

``Cậu nói gì thế?'' Yuka hỏi lại với khuôn mặt trông như không thể tin
được những gì mà tôi vừa nói. ``Lý do tớ bảo cậu tham gia đâu phải để
lấy giải hay là nổi tiếng?''

``Ế? Thế không vì mấy cái đó thì tham gia làm gì?''

``Chán cậu quá, tự nghĩ đi.'' Yuka thở dài, lấy tay chống cằm và nhìn
vào mặt tôi. ``Tớ giận rồi.''

Ơ?

``Tớ xin lỗi, để tớ tự nghĩ vậy.''

Tôi không muốn mất đi người duy nhất chịu đọc truyện của mình đâu.

``...Tất nhiên là dể trau dồi thêm kĩ năng rồi.'' Yuka thở dài đầy chán
nản khi thấy tôi cứ loay hoay với suy nghĩ của mình. ``Có thể truyện của
cậu sẽ không đạt giải, nhưng chắc chắn qua phần thi này thì cậu sẽ có
thể ải thiện lại một thứ gì đó phải chứ?''

Thật sao? Cải thiện được luôn à?

``Mà, dù cậu không đạt giải thì vẫn còn nhiều cơ hội mà.'' Yuka xoa đầu
tôi. ``Dù sao thì tớ sẽ còn ở bên cậu dài dài, nên cứ tự tin là truyện
của cậu sẽ mãi có một độc giả theo dõi đi.''

Yuka là vậy đấy, dù cho thường hay mắng với đánh tôi thì cậu ấy vẫn rất
tốt và luôn ở bên khi tôi cần. Cậu ấy cũng là người duy nhất ủng hộ việc
viết tiểu thuyết này của tôi.

Hửm? Một tiểu thuyết gia kém cỏi, và một cô bạn thân thiết luôn ở bên và
quan tâm chăm sóc à?

``Hình như... tớ có ý tưởng mới rồi.''

*

``Rốt cuộc thì vẫn không được giải nhỉ?'' Yuka nhìn kết quả được công bố
trên trang web của cuộc thi với một vẻ mặt tiếc nuối.

``Vào được top 10 là ngon rồi mà.'' Tôi cười trừ, xoa đầu cậu ấy. ``Với
kĩ năng như tớ thì như vậy đã quá tốt rồi.''

Truyện mà tôi nộp lên để dự thi có thể nói là truyện tốt nhất từ trước
đến giờ bản thân từng viết. Thậm chí là đưa cho bạn học trong câu lạc bộ
xem thì họ cũng khá là trầm trồ. Với lại việc chen chúc qua mấy trăm thí
sinh để đạt được top 10 đã là một thành công quá lớn lao đối với tôi
rồi.

``Nhưng tớ vẫn muốn cậu chiến thắng cơ...'' Yuka tiếp tục nói với tông
giọng buồn bã.

Đúng là fan trung thành nhỉ? Hôm trước vừa bảo không cần lấy giải làm
tôi tưởng rằng khi đạt top 10 thì Yuka sẽ không buồn lắm, ai ngờ lại như
này.

Mà thôi, vì cậu ấy dễ thương nên không sao hết.

``Lần sau thi lại là được mà.'' Tôi lên tiếng trấn an cậu ấy. ``Tớ nhất
định sẽ lấy cái giải nhất cho cậu xem!''

``Cậu vẫn định thi nữa ư?''

``Tất nhiên! Tớ sẽ không bỏ cuộc đâu!''

Dù cho có thất bại bao nhiêu lần, hay sẽ không ai thèm đọc truyện của
tôi đi nữa thì tôi vẫn sẽ tiếp tục viết.

Dù sao thì, tôi vẫn sẽ luôn có một độc giả cho mình mà.

Độc giả, và cũng là người bạn quan trọng nhất của tôi.
}